\color{unimain}{\subsection{Cигнализация}}
\color{black}

\begin{enumerate}

\item
Устройство формирует выходной сигнал на сигнализацию по логической схеме «ИЛИ» при наступлении указанных ниже событий.

\item
Сигналы \textit{«Неисправность»} и \textit{«Неисправность фикс»} формируются при возникновении любого из следующих событий:
\begin{itemize}
\item неисправности в цепях переменного тока ДЗШ;
\item неисправности в цепях переменного напряжения;
\item неисправности в цепях УРОВ СШ;
\item несоответствие положения переключателей в выходных цепях;
\item несоответствие положения испытательных блоков;
\item неисправности в цепях оперативного тока и сигнализации;
\item неисправность терминала;
\item вывод терминала;
\item неселективный режим;
\item Потеря GOOSE.
\end{itemize}

\item
Сигналы \textit{«Срабатывание»} и \textit{«Срабатывание фикс»} формируются при возникновении любого из следующих событий:
\begin{itemize}
\item срабатывание дифференциальной токовой защиты шин;
\item срабатывание УРОВ СШ;
\item прием сигнала внешнего отключения;
\item срабатывание ЛОП.
\end{itemize}

\item
Алгоритм работы логики логики предупредительной сигнализации приведен на рисунке \ref{pic:SIGN}.

\medskip
\begin{figure}[H]
\centering
\includegraphics[width=1\textwidth,height=1\textheight,keepaspectratio]{\fbpath/04.01 ДЗШ/91 Сигнализация/img/SIGN.pdf}
\captionof{figure}{Функциональная схема логики предупредительной сигнализации}\label{pic:SIGN}
\end{figure}

\end{enumerate}